% !TeX root = ../papers/nc-p3.tex
\section{空間三次と四次の七次曲線の排除および七次点グラフの被覆}
\label{sec:reducible-septic-coverage}
\begin{theorem}\label{thm:spatial-cubic-quartic-septic}
点ファイバーを持つ標準AS窓の七次点グラフは、空間有理正規三次 $C$ と整な四次 $D$ の和にはならない。
\end{theorem}
\begin{proof}
七次の純性により点グラフは被約な $Y=C\cup D$ そのものであり $p_a(Y)=5$ である。
平面四次は他の全位置でも算術種数3により平面に入り、隣接平面積の三双一次式の次数上限3に反する。
そこで $D$ は非平面、$h=p_a(D)=0$ または1である。$T=C\cap D$ は長さ $6-h$ で、
$C$ 上では有効Cartier因子である。$h=0$ なら $D$ は滑らかな有理曲線、
$h=1$ なら滑らかな楕円曲線または一点が節点・尖点の種数一曲線で、いずれもGorensteinである。

まず $Y$ に接次元3の点 $v$ がある場合を除く。二つの滑らかな枝の和の接次元は高々2である。
実際、二つのパラメータを交点スキーム $k[[u]]/(u^m)$ 上で一致させると、両局所環のファイバー積は
$k[[u,y]]/(y(y-u^m))$ となる。従って $v$ は $C$ と特異な $D$ の二重特異点との交点に限る。
命題\ref{prop:automatic-spatial-tangent-mark}のOre表示を使う。この点からの水平射影で
$C$ の像は平面二次曲線、$D$ も同様である。後者は正規化上で四次の偏極から二重特異点の
基点次数2を引くことと、非平面性から従う。両者とも各位置でこの二次曲線に双有理に写る。
三生成基底が平面全体の非恒等グラフなら定理\ref{thm:plane-ore-unmarked}により点曲面が入る。
恒等なら持上げの高さを同時に動かせるため同じ矛盾になる。残る三生成基底の点曲線は平面三次であり、
二つの既約な二次曲線を含めず、両水平像は同じである。
そこで一般の水平点で $v(gp)$ と $v(p)$ が独立なら、混合式
$a v(gp)-b v(p)=d_i(v(p),v(gp))$ が両高さを一意に決め、異なる二成分が同じ持上げを占めてしまう。
比例する場合は、持上げのある各点で高さが一次元動き曲面となる。これで空間接点はない。

従って $Y$ は局所平面で、$T$ は $D$ 上でもCartierである。
実際、両成分の共通点では滑らかな局所曲面の中の二つの因子として交わる。
随伴または留数により $\omega_Y|_C=\omega_C(T)$、$\omega_Y|_D=\omega_D(T)$ である。
$L_i$ の多重次数は $(3,4)$ で、$\omega_Y/L_i$ の多重次数は $(1-h,h)$ となる。
Riemann--Rochと四つの非退化な文字から $h^0(\omega_Y/L_i)\ge1$。
次数零側への制限は単射である。そこで零になる切断は、他方の次数1の束の切断として長さ $6-h$ で消えねばならないからである。
次数零の整曲線上の非零切断は生成元となるので $h^0=1$。非零切断は他方でも恒等零でなく、
$T$ と交わらない長さ1のCartier零因子を持つ。従って滑らかな一点 $p_i$ により
$L_i=\omega_Y(-p_i)$ であり、全ての $p_i$ は同じ次数1側成分にある。四文字は完全系となる。

次数1側にのみ台を持ち $T$ を避ける等次数 $e\le4$ の二つの有効因子が線型同値なら、両者は一致する。
その比の有理関数は他方の成分で非零定数であり、前者でも $T$ に沿ってその定数に等しい。
非定数であればこの差は極の総次数 $e$ 以下に対し長さ $6-h>4$ の零点を持ってしまう。
この比較は特異な整四次上でも、$T$ がCartierであり極が滑らかな点に限られるので同じである。
偏極三階差から $p_i+2p_{i+3}\sim2p_{i+1}+p_{i+4}$、従って因子として等しく
$p_{i+3}=p_{i+1}$。三階差に戻すと $4p_i\sim4p_{i+1}$ なのでこれも等しい。
全偏極が等しく、全切断系で座標を揃えれば六関係は $Y$ の対角上で消える。
$Y$ を含む二次式は高々一つである。独立な二式は次数7を切れず、共通平面因子があれば
非平面な $C$ が残る二一次式の直線に入ってしまう。従って対角上に少なくとも二次曲面が残り矛盾する。
\end{proof}

\begin{corollary}\label{cor:all-septic-point-windows}
七次の全標準AS窓は、直線成分がある場合、直線がなく二次曲線成分がある場合、そのいずれもなく
平面三次成分がある場合、および整な七次八係数族で被覆される。
前の三行はそれぞれ系\ref{cor:line-component-atlas}、系\ref{cor:conic-component-atlas}、
定理\ref{thm:plane-cubic-mark-normality}の必要十分な有限条件に純次数7を課す。
行の所属は添字保持同型不変で、印の指定を忘れた比較と最小周期は各行の判定で得られる。
六次で直線・二次曲線・平面三次のどの成分も持たない点グラフが存在すれば、
相異なる二本の有理正規三次の被約な和で、その交わりは長さ4である。
\end{corollary}
\begin{proof}
系\ref{cor:full-quintic-point-window-atlas}の次数分割を使う。
七次で残る $3+4$ の空間三次の場合は前定理で、整七次は定理\ref{thm:integral-septic-boundary}である。
六次は純で、整な場合は定理\ref{thm:integral-sextic-exclusion}、二重三次は前節で除かれた。
二成分の算術種数3から交わりの長さ4が従う。
\end{proof}
